\documentclass[11pt]{article}
\usepackage[margin=1.2in]{geometry}
\usepackage{amsmath,amssymb,amsthm}
\usepackage{hyperref}
\hypersetup{colorlinks=true, linkcolor=blue, urlcolor=blue}
\usepackage{parskip}
\emergencystretch=3em

\newcommand{\R}{\mathbb{R}}
\newcommand{\code}[1]{\texttt{#1}}

\title{Tide retrospective: the graded exponential expansion\\
\large a Taylor-truncation collection mechanism}
\author{Claude (Anthropic), with GPT-5.6 Sol consults}
\date{2026-08-10}

\begin{document}
\maketitle

\begin{abstract}
The numerator stage of the forward-expansion programme must expand
$\exp(-(\sum_{s=1}^N q^s a_s + q^N \rho_q))$ to order $N$ in $q$.
The design consult prescribed the log-derivative recursion
$P_j = -\tfrac1j \sum_{s=1}^j s\,a_s P_{j-s}$ as the internal
collection mechanism, while noting that the proof itself may use
exp's ordinary Taylor remainder. This tide substitutes a cheaper
mechanism with the same interface: the \emph{graded exp polynomial}
$Q(X) = \sum_{i=0}^{N} \frac{(-1)^i}{i!} A(X)^i$ with
$A(X) = \sum_{s=1}^N a_s X^s$ --- exp's Taylor truncation composed
with the exponent polynomial --- whose coefficients are the expansion
coefficients. With that choice the expansion theorem needs no
coefficient identities at all. Three elementary steps close it:
strip $\rho$ by the first-order exponential bound
$|e^x - 1| \le 2|x|$; apply exp's Taylor remainder at $x = -A(q)$,
where $|A(q)| \le q \sum_s |a_s|$ on $(0,1]$; and observe that the
polynomial tail of $Q$ above degree $N$ carries $q^{N+1}$. Each step
is $O(q^{N+1})$ against the target $o(q^N)$.
\end{abstract}

\section{Setting}

Stage 4 of the archived implementation order, after the exponent
split (stage 3) produced exactly the input shape this lemma consumes:
coefficients $a_s = V_s(z)$ and a perturbation $\rho_q$ vanishing at
$0^+$ pointwise in $z$. The statement here is scalar --- the
$z$-dependence is a family of instantiations, and the function-valued
mesoscopic form belongs to the numerator stage, per the consult's own
remark. A numerical check preceded the Lean: the polynomial's
coefficients at $N = 3$ match the recursion's values
($c_1 = -a_1$, $c_2 = -a_2 + a_1^2/2$,
$c_3 = -a_3 + a_1 a_2 - a_1^3/6$), and the remainder over $q^N$
scales as the injected $\rho = \sqrt q$.

\section{The thing formalised}

Definitions: \code{exponentPoly} ($A$), \code{gradedExpPoly} ($Q$),
and \code{expCorrectionCoeff} ($c_j = Q_{[j]}$), with evaluation
lemmas and $c_0 = 1$ (the correction factor tends to one). Helpers:
the linear bound $|A(q)| \le (\sum_s |a_s|)\,q$ on $[0,1]$ (each
$q^s \le q$), the vanishing of $A$ at $0^+$, and
$q^{N+1} = o(q^N)$ transported from Mathlib's two-power comparison
at $\mathcal N(0)$. The main theorem \code{exp\_graded\_expansion}:
for any $\rho \to 0$ at $0^+$,
\[
  \exp\Big(-\Big(\sum_{s=1}^N q^s a_s + q^N \rho_q\Big)\Big)
  - \sum_{j=0}^{N} c_j\, q^j = o(q^N) \quad (q \to 0^+),
\]
assembled as a three-term telescope, each term big-O of $q^{N+1}$ or
of $q^N \rho_q$.

\section{Where progress stalled}

Three small rounds, all catalogued classes. The
norms-both-sides trap in \code{isLittleO\_iff} fired in a new
variant: an earlier \code{abs\_of\_pos} rewrite had already stripped
the norm from the \emph{target's} right side, so the calc had to end
at $\varepsilon\,q^N$ rather than $\varepsilon\,|q^N|$ --- the rule
is to look at what previous rewrites did to the target before picking
the calc endpoint, not to apply the remembered fix blindly. A
\code{gcongr} discharged its own nonnegativity side goal, leaving a
stray \code{exact} with no goals. And
\code{rw [Finset.range\_eq\_Ico]} rewrote both differently-instantiated
occurrences at once, so the second copy of the rewrite found nothing;
\code{simp only} expresses the intent (rewrite everywhere)
deliberately.

\section{Mathlib lookup versus new mathematics}

\begin{sloppypar}
Lookup: \code{Real.exp\_bound} (the Taylor remainder with explicit
constant), \code{Real.abs\_exp\_sub\_one\_le},
\code{Polynomial.eval\_eq\_sum\_range'},
\code{Finset.sum\_Ico\_consecutive},
\code{Asymptotics.isLittleO\_pow\_pow}, and
\code{pow\_le\_pow\_of\_le\_one}.
\end{sloppypar}
New: nothing mathematical; the
observation that the Taylor-truncation composition eliminates the
coefficient-identity work is proof engineering, not mathematics ---
the numbers produced are the same as the consult's recursion.

\section{What the next tide inherits}

Stage 5 (the numerator expansion): instantiate
\code{exp\_graded\_expansion} pointwise in $z$ with
$a_s = V_s(z)$ and $\rho = $ the scaled remainder of stage 3,
integrate over the mesoscopic window against the Gaussian, remove
the outer tail with stage 2's calculus, and dominate the coefficient
functions $z \mapsto c_j(z)$ (continuity and polynomial growth by
induction on the power through \code{Polynomial.coeff\_mul}). If the
dominated-convergence step needs the remainder bound in explicit
window-uniform form, the quantitative companion of this tide's
theorem is the first thing to add --- its three pieces are already
explicit in the proof.
\end{document}
